\section{Conclusions}
\label{sec:conclusion}

In this work, we moved the alert-triage procedure out of small local LLMs and into \hesp{}, a controller that splits an investigation into what to probe next and when the evidence suffices: a hypothesis ledger with counted likelihood tables ranks read-only probes by information gain per cost, an evidence rule and a controller-side stop govern verdicts, and a write-ahead journal makes every step auditable. Four pre-registered studies with five open-weight models and 7{,}272 episodes show that small models fail at the procedure, that choosing probes and deciding to stop are separate failures that different small models exhibit differently, and that taking the procedure away from the model lifts Qwen2.5-7B from 0.125 to 1.000 and Llama-3.1-8B from 0 to 0.917 verified completion. The controller alone, with no LLM, reaches 0.917 as well. In an environment where each cause leaves a distinctive discrete observation, the procedure does the work, and the lesson is about control rather than capability: a small local model should not own the investigation procedure. Whether a model adds value where it matters most, reading ambiguous raw telemetry in which attacks look like routine work and routine work looks like an attack, requires environments and data that this work and the public security datasets we examined do not provide.
